\scanpage{153}
Soit $G$ un groupe algébrique. On désigne par $\deleted{A}K(G)$
l'anneau \deleted{\unsure{des}} engendré par les classes des représentations
rationnelles irréductibles de $G$ : en tant que groupe
additif, il est donc libre ; c'est un \uline{anneau}
\uline{de Chern}. On désigne par $C\deleted{A}K(G)$ ``l'anneau
des classes de Chern'' de $\deleted{A}K(G)$, \inserted{c'est un anneau gradué,} \deleted{$\unsure{A(G)}$ et $\unsure{CA(G)}$ sont des foncteurs}
et on a un homomorphisme canonique
\[
\deleted{A}K(G)\xrightarrow{c_G}\mathbb Z\times(1+C\deleted{A}K(G)^+) \tag*{(1)}
\]
$\deleted{A}K(G)$ et $C\deleted{A}K(G)$ sont des foncteurs en $G$, et
(1) est un hom. fonctoriel. Si $\sigma$ est un
autom. \uline{intérieur} de $G$, alors $\sigma^*$ est l'automorphisme identique de $K(G)$ et $CK(G)$.

Soit $P$ un fibré algébrique principal de
groupe structural $G$, de base $X$. À toute
\deleted{\unsure{quel}} représentation lin. rationnelle $\rho$ de $G$
est associé un fibré vectoriel \inserted{algébrique} sur $X$, d'où
un élément de $\deleted{\unsure{C}}\deleted{A}K(X)\deleted{\unsure{{}_P}}$. On en
conclut un hom. d'anneaux de Chern
\[
K(G)\longrightarrow K(X)
\]
d'où un hom d'anneaux gradués
\[
CK(G)\longrightarrow CK(X)=A(X) \tag*{(2)}
\]
appelé \uline{hom. caractéristique} du fibré $P$. On
a commutativité des
% archive sheet number: 152
% printed form (vertical paper watermark)
\unsure{R|ER} \tikz[baseline=-.4ex,scale=.15]{\foreach \a in {0,45,...,315}{\draw (\a:.52) circle (.22);} \draw (0,0) circle (.21);} ARGENTINE \tikz[baseline=-.2ex,scale=.17]{\draw (0,0)--(1.6,0)--(1.8,.33)--(.5,.33)--(.32,.13)--(.92,.13);} TYPEWRITER \tikz[baseline=-.4ex,scale=.15]{\foreach \a in {0,45,...,315}{\draw (\a:.52) circle (.22);} \draw (0,0) circle (.21);} \unsure{A}
